\chapter{Inventory Before Provisioning}

Provisioning instructions become dangerous when they assume that every host is
empty and identically configured. Before choosing addresses or jail names, we
will collect a small, read-only inventory. The result tells us what BSDRunner
can reuse and where a new service could collide with existing work.

This chapter does not create, start, stop, or modify a jail.

\section{Exercise 2: Inspect the host}

From the BSDRunner repository on the FreeBSD host, inspect the operating system,
ZFS storage, and resource-control support:

\begin{lstlisting}[language=sh]
freebsd-version -k
freebsd-version -r
freebsd-version -u
zpool list -o name,size,alloc,free,health
sysctl kern.racct.enable
\end{lstlisting}

The three release lines report the installed kernel, running kernel, and
installed userland, in that order. They will often agree. If the installed and
running kernel patch levels differ, an update may be waiting for a reboot. Note
the difference; do not turn an inventory exercise into an unplanned reboot.

Record the pool we will use only if its health is \texttt{ONLINE} and it has
enough free space for four small jails, packages, gems, and PostgreSQL data.

A value of \texttt{1} for \texttt{kern.racct.enable} means that Blackship can
enforce RCTL limits. RCTL is useful, but it is not a prerequisite for the Rails,
nginx, or PostgreSQL starters in this tutorial. If the value is \texttt{0}, note
it rather than changing a boot setting now.

\section{Inspect Blackship's contract}

BSDRunner requires Blackship 0.6.2 or newer and machine-output schema version 1.
The capability document also shows which FreeBSD facilities Blackship can see:

\begin{lstlisting}[language=sh]
blackship --format json capabilities | jq '.data | {
  blackship_version,
  output_schema_version,
  platform: {
    os: .platform.os,
    arch: .platform.arch,
    kernel_release: .platform.kernel_release
  },
  facilities: [.facilities[] | {name, available}]
}'
\end{lstlisting}

Look for an available \texttt{jail}, \texttt{jls}, \texttt{jexec}, \texttt{pf},
and \texttt{zfs}. This is a capability check, not a request to set up or change
any of them.

\section{Ask BSDRunner for the consolidated view}

The individual commands establish what the host can do. BSDRunner's snapshot
adds project state, current service names, and the managed private network:

\begin{lstlisting}[language=sh]
sh dotfiles/.config/bsdrunner/scripts/\
bsdrunner-jails-backend.sh snapshot | jq '{
  host_release,
  storage: (.storage | {pools, zfs_pool}),
  running,
  configured,
  blackship: (.blackship | {
    available,
    compatible,
    version,
    configured,
    recommended_address,
    rctl_enabled,
    networks: .networks.data,
    managed
  })
}'
\end{lstlisting}

In this output, \texttt{blackship.configured} means that BSDRunner's managed
project manifest exists. An active \texttt{bsdrunner} network gives the four new
jails a private subnet. The recommended address is only the next candidate; no
address has been reserved by running this command.

The top-level \texttt{running} and \texttt{configured} arrays include native
and Bastille jails. The \texttt{managed} array inside \texttt{blackship} is the
important collision check for BSDRunner-managed services. Existing services
are welcome, but none may already use one of our proposed names:

\begin{lstlisting}
web
store
catalog
database
\end{lstlisting}

After a reboot, the network may be listed with status \texttt{missing-bridge}.
Check whether BSDRunner's Blackship boot service is enabled:

\begin{lstlisting}[language=sh]
sysrc bsdrunner_blackship_enable
\end{lstlisting}

When the result is \texttt{NO} because every managed jail has autostart disabled,
\texttt{missing-bridge} is expected: the durable network definition exists, but
its runtime interface has not been created. A guarded BSDRunner start restores
Blackship networking before it starts the requested jail. Do not create the
bridge by hand. An absent \texttt{bsdrunner} network entry is a different state
and does block provisioning.

\section{Checkpoint}

Before moving on, record these facts:

\begin{description}
  \item[FreeBSD releases] the installed kernel, running kernel, and userland;
  \item[ZFS pool] an online pool with sufficient free space;
  \item[Blackship] version 0.6.2 or newer, schema 1, and compatible;
  \item[private network] its name, subnet, gateway, bridge, and runtime state;
  \item[next address] BSDRunner's recommended unused address; and
  \item[name check] whether \texttt{web}, \texttt{store}, \texttt{catalog}, and
    \texttt{database} are all unused.
\end{description}

Stop here if Blackship is incompatible, the network definition is absent, the
pool is not healthy, a proposed name is already managed, or the kernel and
userland differ by more than a patch level. A missing runtime bridge is also a
reason to stop when the Blackship boot service is enabled; with autostart
disabled, it is an expected idle state. These are planning inputs, not errors to
work around blindly.

Once the checkpoint is satisfied, the next exercise will create only the
PostgreSQL jail. We will verify its lifecycle and health before introducing a
Rails application.
